% GENERATED FILE. Edit canonical lessons, then run npm run generate:books.
% canonical-source-hash: fnv1a64:e0af9012985539c7
% canonical-lessons: SA-C529-parivartanam, SA-C529-upasamharati, SA-C529-samsodhayati, SA-C529-tvarayati, SA-C529-tvara

\chapter{A Change of Plan}
\label{ch:sa-a-change-of-plan}
\hlchaptermodality{529}

\begin{chapteropening}
\textbf{By the end of this chapter:} \emph{I can say what changes if something goes wrong.}

Complete the last lesson of chapter 529: I can say what changes if something goes wrong.
\end{chapteropening}

\section[\texorpdfstring{\sk{परिवर्तनम्} (parivartanam)}{परिवर्तनम् (parivartanam)}]{\sk{परिवर्तनम्} (parivartanam) — a change; an exchange}

\label{lesson:SA-C529-parivartanam}



\begin{quote}
Before the new one: say the Sanskrit for makes ready, then the Sanskrit for an almanac.
\end{quote}

\subsection*{You'll want to know: \sk{परिवर्तनम्}}
\textbf{\sk{परिवर्तनम्}} — \emph{parivartanam} — \textquotedblleft{}a change; an exchange\textquotedblright{}.

\begin{quote}
\sk{यदि} \sk{वर्षा} \sk{भवति}, \sk{तर्हि} \sk{योजनायां} \sk{परिवर्तनं} \sk{भविष्यति।}
\end{quote}

\begin{quote}
\emph{yadi varṣā bhavati, tarhi yojanāyāṃ parivartanaṃ bhaviṣyati.}
\end{quote}

\begin{quote}
If it rains, the plan will change.
\end{quote}

A neuter noun.

\begin{cousinweb}
\textbf{\sk{परिवर्तनम्}} (\emph{parivartanam}) is built on \emph{pari}, \textquotedblleft{}around\textquotedblright{}, and the root √vṛt, \textquotedblleft{}to turn\textquotedblright{}, which also gives \emph{parivartayati}; MW gives \emph{parivartana} as turning round, and barter, exchange; \textquotedblleft{}a change of plan\textquotedblright{} grows from these senses (MW).
\end{cousinweb}

\begin{grammarlens}[title={What you've built}]
One more word for saying what happens if.
\end{grammarlens}

\subsection*{Your turn}
\begin{itemize}
\raggedright
  \item \emph{Say it:} \emph{parivartanam}
  \item \emph{Say it:} \emph{parivartanam}, once more
  \item \emph{Say it:} \emph{yadi varṣā bhavati, tarhi yojanāyāṃ parivartanaṃ bhaviṣyati.} — If it rains, the plan will change.
  \item \emph{Recall:} say the Sanskrit for makes ready, then the Sanskrit for an almanac, then say \emph{parivartanam} again
\end{itemize}

\begin{tcolorbox}[breakable,colback=teal!4,colframe=teal!35!black,title={Before you move on}]
What does \sk{परिवर्तनम्} mean? (\textquotedblleft{}A change; an exchange\textquotedblright{}.) Say it once more.
\end{tcolorbox}
\section[\texorpdfstring{\sk{उपसंहरति} (upasaṃharati)}{उपसंहरति (upasaṃharati)}]{\sk{उपसंहरति} (upasaṃharati) — winds up; withdraws}

\label{lesson:SA-C529-upasamharati}



\begin{quote}
Before the new one: say the Sanskrit for an almanac, then the Sanskrit for a change.
\end{quote}

\subsection*{You'll want to know: \sk{उपसंहरति}}
\textbf{\sk{उपसंहरति}} — \emph{upasaṃharati} — \textquotedblleft{}winds up; withdraws\textquotedblright{}.

\begin{quote}
\sk{यदि} \sk{अतिथिः} \sk{न} \sk{आगच्छति}, \sk{तर्हि} \sk{माता} \sk{उत्सवम्} \sk{उपसंहरति।}
\end{quote}

\begin{quote}
\emph{yadi atithiḥ na āgacchati, tarhi mātā utsavam upasaṃharati.}
\end{quote}

\begin{quote}
If the guest does not come, my mother winds up the party.
\end{quote}

Past (he …): \sk{सः} \sk{सभाम्} \sk{उपसंहृतवान्}, \emph{saḥ sabhām upasaṃhṛtavān}.

\begin{cousinweb}
\textbf{\sk{उपसंहरति}} (\emph{upasaṃharati}) is built on \emph{upa}, \emph{sam}, \textquotedblleft{}together\textquotedblright{}, and the root √hṛ, \textquotedblleft{}to carry\textquotedblright{}, which also gives \emph{harati}; MW gives \emph{upasaṃhṛ} as to draw together, to sum up, to withdraw, to stop (MW).
\end{cousinweb}

\begin{grammarlens}[title={What you've built}]
One more word for saying what happens if.
\end{grammarlens}

\subsection*{Your turn}
\begin{itemize}
\raggedright
  \item \emph{Say it:} \emph{upasaṃharati}
  \item \emph{Say it:} \emph{upasaṃharati}, once more
  \item \emph{Say it:} \emph{yadi atithiḥ na āgacchati, tarhi mātā utsavam upasaṃharati.} — If the guest does not come, my mother winds up the party.
  \item \emph{Recall:} say the Sanskrit for an almanac, then the Sanskrit for a change, then say \emph{upasaṃharati} again
\end{itemize}

\begin{tcolorbox}[breakable,colback=teal!4,colframe=teal!35!black,title={Before you move on}]
What does \sk{उपसंहरति} mean? (\textquotedblleft{}Winds up; withdraws\textquotedblright{}.) Say it once more.
\end{tcolorbox}
\section[\texorpdfstring{\sk{संशोधयति} (saṃśodhayati)}{संशोधयति (saṃśodhayati)}]{\sk{संशोधयति} (saṃśodhayati) — revises; corrects}

\label{lesson:SA-C529-samsodhayati}



\begin{quote}
Before the new one: say the Sanskrit for a change, then the Sanskrit for winds up.
\end{quote}

\subsection*{You'll want to know: \sk{संशोधयति}}
\textbf{\sk{संशोधयति}} — \emph{saṃśodhayati} — \textquotedblleft{}revises; corrects\textquotedblright{}.

\begin{quote}
\sk{यदि} \sk{योजना} \sk{कठिना} \sk{अस्ति}, \sk{तर्हि} \sk{वयं} \sk{तां} \sk{संशोधयामः।}
\end{quote}

\begin{quote}
\emph{yadi yojanā kaṭhinā asti, tarhi vayaṃ tāṃ saṃśodhayāmaḥ.}
\end{quote}

\begin{quote}
If the plan is hard, we revise it.
\end{quote}

Past (he …): \sk{सः} \sk{योजनां} \sk{संशोधितवान्}, \emph{saḥ yojanāṃ saṃśodhitavān}.

\begin{cousinweb}
\textbf{\sk{संशोधयति}} (\emph{saṃśodhayati}) is built on \emph{sam}, \textquotedblleft{}thoroughly\textquotedblright{}, and the causative of the root √śudh, \textquotedblleft{}to be pure\textquotedblright{}, which also gives \emph{śodhayati}; MW gives \emph{saṃśodhayati} as to cleanse thoroughly, to clear; \textquotedblleft{}to revise\textquotedblright{} is its modern use (MW).
\end{cousinweb}

\begin{grammarlens}[title={What you've built}]
One more word for saying what happens if.
\end{grammarlens}

\subsection*{Your turn}
\begin{itemize}
\raggedright
  \item \emph{Say it:} \emph{saṃśodhayati}
  \item \emph{Say it:} \emph{saṃśodhayati}, once more
  \item \emph{Say it:} \emph{yadi yojanā kaṭhinā asti, tarhi vayaṃ tāṃ saṃśodhayāmaḥ.} — If the plan is hard, we revise it.
  \item \emph{Recall:} say the Sanskrit for a change, then the Sanskrit for winds up, then say \emph{saṃśodhayati} again
\end{itemize}

\begin{tcolorbox}[breakable,colback=teal!4,colframe=teal!35!black,title={Before you move on}]
What does \sk{संशोधयति} mean? (\textquotedblleft{}Revises; corrects\textquotedblright{}.) Say it once more.
\end{tcolorbox}
\section[\texorpdfstring{\sk{त्वरयति} (tvarayati)}{त्वरयति (tvarayati)}]{\sk{त्वरयति} (tvarayati) — hurries along}

\label{lesson:SA-C529-tvarayati}



\begin{quote}
Before the new one: say the Sanskrit for winds up, then the Sanskrit for revises.
\end{quote}

\subsection*{You'll want to know: \sk{त्वरयति}}
\textbf{\sk{त्वरयति}} — \emph{tvarayati} — \textquotedblleft{}hurries along\textquotedblright{}.

\begin{quote}
\sk{यदि} \sk{समयः} \sk{अल्पः} \sk{अस्ति}, \sk{तर्हि} \sk{पिता} \sk{अस्मान्} \sk{त्वरयति।}
\end{quote}

\begin{quote}
\emph{yadi samayaḥ alpaḥ asti, tarhi pitā asmān tvarayati.}
\end{quote}

\begin{quote}
If time is short, my father hurries us along.
\end{quote}

Past (he …): \sk{सः} \sk{अस्मान्} \sk{त्वरितवान्}, \emph{saḥ asmān tvaritavān}.

\begin{cousinweb}
\textbf{\sk{त्वरयति}} (\emph{tvarayati}) comes from the causative of the root √tvar, \textquotedblleft{}to hurry, to make haste\textquotedblright{}; MW gives \emph{tvarayati} as to cause to hasten, to urge forward (MW).
\end{cousinweb}

\begin{grammarlens}[title={What you've built}]
One more word for saying what happens if.
\end{grammarlens}

\subsection*{Your turn}
\begin{itemize}
\raggedright
  \item \emph{Say it:} \emph{tvarayati}
  \item \emph{Say it:} \emph{tvarayati}, once more
  \item \emph{Say it:} \emph{yadi samayaḥ alpaḥ asti, tarhi pitā asmān tvarayati.} — If time is short, my father hurries us along.
  \item \emph{Recall:} say the Sanskrit for winds up, then the Sanskrit for revises, then say \emph{tvarayati} again
\end{itemize}

\begin{tcolorbox}[breakable,colback=teal!4,colframe=teal!35!black,title={Before you move on}]
What does \sk{त्वरयति} mean? (\textquotedblleft{}Hurries along\textquotedblright{}.) Say it once more.
\end{tcolorbox}
\section[\texorpdfstring{\sk{त्वरा} (tvarā)}{त्वरा (tvarā)}]{\sk{त्वरा} (tvarā) — haste}

\label{lesson:SA-C529-tvara}



\begin{quote}
Before the new one: say the Sanskrit for revises, then the Sanskrit for hurries along.
\end{quote}

\subsection*{You'll want to know: \sk{त्वरा}}
\textbf{\sk{त्वरा}} — \emph{tvarā} — \textquotedblleft{}haste\textquotedblright{}.

\begin{quote}
\sk{यदि} \sk{त्वरा} \sk{न} \sk{अस्ति}, \sk{तर्हि} \sk{मन्दं} \sk{गच्छतु।}
\end{quote}

\begin{quote}
\emph{yadi tvarā na asti, tarhi mandaṃ gacchatu.}
\end{quote}

\begin{quote}
If there is no hurry, go slowly.
\end{quote}

A feminine noun.

\begin{cousinweb}
\textbf{\sk{त्वरा}} (\emph{tvarā}) comes from the root √tvar, \textquotedblleft{}to hurry\textquotedblright{}; MW gives \emph{tvarā} as haste, speed (MW).
\end{cousinweb}

\begin{grammarlens}[title={What you've built}]
One more word for saying what happens if.
\end{grammarlens}

\subsection*{Your turn}
\begin{itemize}
\raggedright
  \item \emph{Say it:} \emph{tvarā}
  \item \emph{Say it:} \emph{tvarā}, once more
  \item \emph{Say it:} \emph{yadi tvarā na asti, tarhi mandaṃ gacchatu.} — If there is no hurry, go slowly.
  \item \emph{Recall:} say the Sanskrit for revises, then the Sanskrit for hurries along, then say \emph{tvarā} again
\end{itemize}

\begin{tcolorbox}[breakable,colback=teal!4,colframe=teal!35!black,title={Before you move on}]
What does \sk{त्वरा} mean? (\textquotedblleft{}Haste\textquotedblright{}.) Say it once more.
\end{tcolorbox}
